\documentclass[11pt,a4paper]{article}
\usepackage[margin=25mm]{geometry}
\usepackage[T1]{fontenc}
\usepackage{lmodern,amsmath,amssymb,amsthm,booktabs,longtable}
\usepackage[hidelinks]{hyperref}
\usepackage{microtype}
\newcommand{\Q}{\mathbb Q}
\newcommand{\F}{\mathbb F}
\newcommand{\Z}{\mathbb Z}
\newcommand{\Pic}{\operatorname{Pic}}
\newcommand{\Jac}{\operatorname{Jac}}
\newcommand{\A}{\mathcal A}
\newcommand{\W}{W^0_5}
\newtheorem{theorem}{Theorem}
\newtheorem{lemma}[theorem]{Lemma}
\title{Quintic points on thirty remaining modular curves\\
  \large A consolidated proof and reproducible certificates}
\author{}
\date{5 September 2026}
\begin{document}
\maketitle
\begin{abstract}
We prove finiteness of degree-five closed points over $\Q$ for the thirty
levels left open in Section~7 of Derickx--Hwang--Jeon--Orli\'c
\cite{DHJO}. None admits a degree-five function over $\Q$.
The proof uses exhaustive finite-field divisor calculations and one local
lifting lemma. This note separates the mathematical deductions from the
computations. All necessary Magma code and verification output accompany it.
\end{abstract}

\section{Result and reading guide}
\begin{theorem}\label{thm:main}
For every level in Table~\ref{tab:overview}, $X_0(N)$ has only finitely many
closed points of degree exactly five over $\Q$, and has no function of
degree exactly five over $\Q$.
\end{theorem}
Finiteness does not mean emptiness: we do not enumerate the quintic points.
A degree-five point means one closed point with residue field of degree
five, not an arbitrary effective divisor of degree five.
No step assumes BSD or equality between analytic and algebraic rank.

There are two computational ideas:
\emph{enumerate all degree-five divisor classes over a small finite field},
and \emph{lift a surviving family when its tangent space has the expected
dimension}. Three routes organize their application:
\begin{itemize}
\item R: direct reduction (rank zero), or no surviving subgroup coset;
\item Q: a rank-zero quotient labels all possible abelian translates;
\item T: tangent spaces force surviving translates into a known family.
\end{itemize}
The function column records the separate obstruction to a degree-five
function. The double-cover argument is a one-line genus inequality.
Local pencils use the same lifting lemma as route T.

\begin{longtable}{@{}rrll@{}}
\caption{All thirty cases and their proofs.}\label{tab:overview}\\
\toprule $N$&$g$&Finiteness route&Degree-five function obstruction\\
\midrule\endfirsthead
\toprule $N$&$g$&Finiteness route&Degree-five function obstruction\\
\midrule\endhead
74&8&T&Local pencil, two-prime sieve\\
77&7&T&Local pencil\\
79&6&Q&Quotient sieve\\
82&9&T&Published $\Q$-gonality $6$\\
83&7&Q&Double cover\\
85&7&R (empty cosets)&Local pencil\\
86&10&T&$W^1_5(\F_7)=\varnothing$\\
88&9&R (empty cosets)&Local pencil\\
89&7&Q&Double cover\\
92&10&R (empty cosets)&Double cover\\
96&9&R (rank zero)&Rank-zero lifting\\
98&7&R (rank zero)&Rank-zero lifting, with $X_0(49)$\\
99&9&T&$W^1_5(\F_5)=\varnothing$\\
100&7&R (rank zero)&Rank-zero lifting\\
101&8&Q&Double cover\\
103&8&Q&Quotient sieve\\
107&9&Q&Double cover\\
111&11&T&Double cover\\
118&14&T&Double cover\\
123&13&T&Double cover\\
128&9&T&$W^1_5(\F_3)=\varnothing$\\
131&11&Q&Double cover\\
141&15&T&Double cover\\
142&17&T&Double cover\\
143&13&T&Double cover\\
145&13&T&Double cover\\
155&15&T&Double cover\\
159&17&T&Double cover\\
167&14&Q&Double cover\\
191&16&Q&Double cover\\
\bottomrule
\end{longtable}

\section{The common framework}\label{sec:framework}
Put $C=X_0(N)$ and $J=\Jac(C)$. A rational cusp identifies $\Pic^5(C)$
with $J$. Write $W^r_d$ for the locus of degree-$d$ line bundles $L$
with $h^0(L)\ge r+1$. Since $C$ has a rational point, rational Picard
classes are represented by line bundles over $\Q$; there is no Brauer
obstruction.

Suppose first that $C$ has no degree-five function over $\Q$.
An irreducible rational effective divisor $D$ of degree five must have
$h^0(D)=1$: a nonconstant rational section would have a nonempty
rational pole divisor bounded by $D$, necessarily equal to $D$.
Thus quintic points inject into $\W(\Q)$.

Factor $J$ up to $\Q$-isogeny into modular simple factors. Let $\A$
be the sum of the isotypic parts whose central $L$-value vanishes.
Exact modular-symbol central-value tests identify these parts.
The nonvanishing theorem for modular abelian varieties \cite{KL}
gives $\operatorname{rank}(J/\A)(\Q)=0$.
We use only this implication: vanishing factors are \emph{possible}
positive-rank factors, without claiming their exact ranks.

Mordell--Lang \cite{Faltings} says that the Zariski closure of $\W(\Q)$
is a finite union of points and abelian translates. Every positive-dimensional
component with dense rational points is defined over $\Q$ and its direction
is contained in $\A$, since its image in the rank-zero quotient is finite.
It suffices to show that all such components have fixed proper
sub-divisors on their $h^0=1$ loci. The remaining finitely many classes
then give finitely many quintic points.

\subsection{What the finite-field enumeration certifies}
At each working prime $p\nmid N$, the scripts check that the reduced
canonical model is integral, has a rational place, and that its arithmetic
genus equals its function-field genus. It is therefore smooth and
geometrically integral and represents the good reduction of $C$.

Every effective degree-five divisor over $\F_p$ is a sum of closed places
of degrees at most five. For each partition of five, we enumerate the
required multisets of places. This is exhaustive, with no height bound.
The class group calculation gives all of $\W(\F_p)$: over a finite field
every rational Picard class is represented by a line bundle.

If $x+B\subset\W$ over $\Q$, its reduction contains the entire coset
$\bar x+B(\F_p)$ in $\W(\F_p)$, not just reductions of its rational
points. This follows from geometric containment and the smooth abelian
models at good reduction. One missing element excludes the translate.

\subsection{One local lifting lemma}\label{sec:local}
\begin{lemma}[Local rigidity]\label{lem:local}
Let $R=\Z_p$, let $Z$ be a closed subscheme of a scheme smooth over $R$,
and let $U\to Z$ be a morphism with $U$ smooth over $R$ of relative
dimension $n$. At rational special-fiber points $u\mapsto z$, suppose
\[
 \dim_{\F_p}T_z Z_{\F_p}=n,\qquad
 T_u U_{\F_p}\longrightarrow T_z Z_{\F_p}\text{ is an isomorphism}.
\]
Then the morphism identifies the completed local rings.
\end{lemma}
\begin{proof}
Lifts of relative cotangent generators give a surjection
$R[[x_1,\ldots,x_n]]\twoheadrightarrow\widehat{\mathcal O}_{Z,z}$.
Its composite to $\widehat{\mathcal O}_{U,u}\simeq R[[t_1,\ldots,t_n]]$
has invertible linear part modulo $p$, so is an isomorphism by formal
inversion. The surjection has zero kernel. Flatness of $Z$ is not assumed.
\end{proof}

We apply this to $Z_S=\bigcap_{s\in S}(\W-s)$.
Prime-to-$p$ rational torsion points on the special fiber lift uniquely
to torsion over $R$. If $S$ lies in both a candidate direction and a
known $B$, the same lifts belong to both.
When the direction is already known to equal $B$, points of $B(\F_p)$
can instead be lifted arbitrarily inside $B(R)$, including $p$-primary
points: both families are stabilized by all of $B$.

At a class with $h^0(D)=1$, the annihilator of $T_{[D]}\W$ in $H^0(K_C)$
is $H^0(K_C-D)$. For $h^0(D)\ge2$, the tangent space is the full ambient
space. These are the first-order determinantal equations of $W^0_5$.
The tangent space of $Z_S$ is computed by spanning these annihilators
over the coset, in one fixed basis of $H^0(K_C)$.

If the lemma identifies the local branch with a known proper image $V$,
a candidate translate through that residue class is locally contained
in $V$. Properness makes $V$ closed. A nonempty $p$-adic open subset
of the smooth translate is Zariski dense, so the entire translate lies in $V$.

\section{Excluding degree-five functions}\label{sec:functions}
\subsection{The double-cover inequality: seventeen cases}
A degree-two map from $C$ to a genus-$h$ curve and a degree-five map
to $\mathbb P^1$ generate the function field, since their degrees are
coprime. Castelnuovo--Severi gives $g(C)\le2h+4$.
Every row below violates this inequality, excluding degree-five functions
even geometrically. The script \texttt{arithmetic.m} checks the canonical
involutions and their invariant dimensions.
\begin{center}\small
\begin{tabular}{@{}rrrr@{\qquad}rrrr@{}}
\toprule
$N$&$g$&involution&$h$&$N$&$g$&involution&$h$\\
\midrule
83&7&$w_{83}$&1&89&7&$w_{89}$&1\\
92&10&$w_{23}$&1&101&8&$w_{101}$&1\\
107&9&$w_{107}$&2&111&11&$w_{111}$&2\\
118&14&$w_{59}$&3&123&13&$w_{41}$&3\\
131&11&$w_{131}$&1&141&15&$w_{47}$&3\\
142&17&$w_{71}$&2&143&13&$w_{143}$&2\\
145&13&$w_{29}$&4&155&15&$w_{155}$&4\\
159&17&$w_{159}$&4&167&14&$w_{167}$&2\\
191&16&$w_{191}$&2&&&&\\
\bottomrule
\end{tabular}
\end{center}

\subsection{Empty special loci and local pencils}
The $\Q$-gonality of $X_0(82)$ is six \cite[Theorem~1.4]{NO}.
At $86,99,128$, respectively at $p=7,5,3$, the numbers of effective
degree-five divisors and their distinct classes are equal. Every class
has $h^0=1$, so $W^1_5(\F_p)$ is empty. Upper semicontinuity excludes
a rational degree-five line bundle with two sections.
These checks are part of \texttt{translates.m}.

For $77,85,88$, \texttt{local\_pencils.m} checks that every special
degree-five divisor has $h^0=2$, moving degree four, and belongs to
$A_4+C$ for one degree-four pencil class $A_4$.
\begin{center}
\begin{tabular}{@{}rrrrr@{}}
\toprule $N$&$p$&effective divisors&special divisors&$\dim T W^1_5$\\
\midrule
77&5&12528&48&1\\
85&7&28676&64&1\\
88&3&3928&32&1\\
\bottomrule
\end{tabular}
\end{center}
The tangent dimension follows from the basepoint-free pencil trick:
\[
                 \dim T_{[D]}W^1_5=8-g+h^0(K_C-D-A_4).
\]
The Petri map has source dimension $2(g-4)$ and kernel
$H^0(K_C-D-A_4)$, which proves the formula.

These curves have $\Q$-gonality four \cite[Theorem~1.2]{NO}.
A rational tetragonal pencil lifts the unique class $A_4$. Its reduction
cannot acquire an extra section, since adding a rational point would
contradict the checked degree-five bound $h^0\le2$.
The Abel map $C\to W^1_5$, $P\mapsto A_4+P$, is an immersion.
Lemma~\ref{lem:local} identifies every relevant formal neighborhood
with this family. Its two sections come from $A_4$, so the added point
is fixed. A basepoint-free rational degree-five pencil is impossible.

Level $74$ uses the same argument plus a two-prime sieve.
Its four rational cusps generate a group of order $513$.
Here cusp differences are torsion by Manin--Drinfeld, and injectivity
at an odd good prime identifies their rational group with its reduction.
The quotient by the possible positive part $E_{37}^2$ has rational
torsion order dividing $513$. At $p=3,5$, the cusp reductions still
generate a group of order $513$, disjoint from the respective primary
groups below. The cusp group therefore maps isomorphically to the
rational rank-zero quotient and supplies common labels.
\begin{center}
\begin{tabular}{@{}rrrr@{}}
\toprule $p$&primary group used&order&exceptional cusp labels\\
\midrule
3&$J(\F_3)[7^\infty]$&49&$(-2,9,0),\ (0,-2,-2)$\\
5&$J(\F_5)[2^\infty]$&256&$(0,-2,1),\ (1,6,0)$\\
\bottomrule
\end{tabular}
\end{center}
At each prime there is one moving tetragonal class. Outside the two
exceptional labels, all eligible special classes have moving degree four
and tangent dimension one. The exceptional sets are disjoint
\emph{in the same cusp group}; the script checks shared integral
representatives, not just inequality of triples.
Every rational class is thus covered by local rigidity at one prime.
The rational tetragonal pencil again comes from \cite{NO}.
The file \texttt{pencils74.m} includes the arithmetic, cusp identification,
and both sieves.

The remaining function exclusions appear below: rank-zero lifting for
$96,98,100$ and the quotient sieve for $79,103$.

\section{Route R: direct reduction}\label{sec:direct}
\subsection{Rank zero and the argument for \texorpdfstring{$X_1(32)$}{X1(32)}}
This is where Section~5.5, Proposition~5.5 of Najman--Varivoda
\cite{NV} applies. All three Jacobians at $96,98,100$ have nonzero
central $L$-value, hence rank zero. Reduction on $J(\Q)$ is injective
at the odd good primes below.
The script \texttt{rank\_zero.m} enumerates every special degree-five
divisor, checks $h^0=2$ and moving degree four, and explicitly lifts
every moving pencil class by a divisor supported on rational points.
\begin{center}\small
\begin{tabular}{@{}rrrrrr@{}}
\toprule
$N$&$p$&$\#C(\F_p)$&rational lifts&degree-four pencil classes&
$\#W^1_5(\F_p)$\\
\midrule
96&7&8&8&2&16\\
98&3&6&4&1&6\\
100&3&6&6&3&18\\
\bottomrule
\end{tabular}
\end{center}
For $96,100$, write the reduction of a rational special class as
$\bar A_4+\bar P$. Lift both terms to $\Q$. Injectivity of reduction
forces the original class to be $A_4+P$. Its two sections come from
$A_4$, so $P$ is fixed and no degree-five function results.
Thus the proposition applies directly. Rank zero gives finitely many
degree-five classes; without degree-five functions each has at most one
irreducible quintic representative.

\paragraph{The necessary refinement at $98$.}
The point-lifting hypothesis fails: six finite-field points versus four
rational lifts. The proposition cannot be applied unchanged.
Use the degree-three degeneracy map
$\pi:X_0(98)\to X_0(49)=E$.
The code verifies $E(\Q)\simeq\Z/2\Z$ and
$E(\F_3)\simeq\Z/4\Z$.
Both rational points of $E$ are its rational cusps.
The other cusps of $E$ form an orbit of degree six; on $X_0(98)$
there are four rational cusps and two degree-six cusp orbits.
The cyclotomic cusp-field description and the order six of $3$ modulo
$7$ show that these other cusps are not $\F_3$-rational.
The four supplied points on $X_0(98)$ are checked to be its
Atkin--Lehner cusp orbit.

The other two points $\bar P\in X_0(98)(\F_3)$ are noncuspidal.
Degeneracy maps take precisely cusps to cusps, also at this good prime,
so $\pi(\bar P)$ is a noncuspidal point of $E(\F_3)$ of order four.
It is not the reduction of a point of $E(\Q)$.
In particular $\bar P$ cannot lift even as a \emph{rational degree-one
Picard class}: pushing such a class forward would give a point of
$\Pic^1(E)(\Q)=E(\Q)$ reducing to $\pi(\bar P)$.

Lift the unique degree-four pencil $A_4$. For a rational special
degree-five class $L$, the rational class $L-A_4$ has degree one.
Its reduction must therefore be a rational cusp, and injectivity gives
$L=A_4+P$ with a fixed rational cusp. This proves both assertions at $98$.

\subsection{Empty cosets: 85, 88, 92}
Each has a unique possible positive-rank simple factor, of multiplicity one.
\begin{center}
\begin{tabular}{@{}rrrrr@{}}
\toprule $N$&$\dim\A$&$p$&$\#\A(\F_p)$&subgroup order\\
\midrule
85&2&3&34&17\\
88&1&3&7&7\\
92&1&3&7&7\\
\bottomrule
\end{tabular}
\end{center}
The indicated prime-primary subgroup of $J(\F_3)$ has that prime order
and belongs to $\A(\F_3)$; the complementary Frobenius order is coprime
to it. The exhaustive test finds no whole coset inside $\W(\F_3)$.
Any positive-dimensional component would have direction $\A$ and
would contain such a coset on reduction. Hence none exists.
The function exclusions already proved complete the finiteness argument.
The file is \texttt{empty\_cosets.m}.

\section{Route Q: nine prime levels}\label{sec:quotient}
Let $Y=C/\langle w_N\rangle$, $\pi:C\to Y$,
$A=\pi^*\Jac(Y)$, and $R=J/A$.
Let $P$ be the infinity cusp, $Q=w_N(P)$, $t=[Q-P]$, and $F=P+Q$.
The cover is ramified, so pullback on Jacobians is injective.
In each row below, $A$ is $\Q$-simple of multiplicity one, and all
other factors have nonzero central $L$-value. The script identifies it
by dimension and, at $89$, by Frobenius order to distinguish another
elliptic factor. Thus $R$ has rank zero.
\begin{center}\small
\begin{tabular}{@{}rrrrrrl@{}}
\toprule $N$&$g$&$g(Y)$&$p$&$m$&$\#A(\F_p)$&surviving indices\\
\midrule
79&6&1&3&13&5&$1,2,3,4$\\
83&7&1&7&41&11&$1,2,3,4,23$\\
89&7&1&3&22&5&$1,2,3,4$\\
101&8&1&5&25&7&$1,2,3,4$\\
103&8&2&5&17&55&$2,3$\\
107&9&2&3&53&29&$2,3$\\
131&11&1&3&65&5&$1,2,3,4$\\
167&14&2&3&83&19&$2,3$\\
191&16&2&2&95&11&$2,3$\\
\bottomrule
\end{tabular}
\end{center}
The gcd of $\#R(\F_q)$ over odd primes $q=3,5,7,11,13,17,19$
is $m$. Rational torsion injects at these primes, so $\#R(\Q)$ divides $m$.
At the working prime, the code constructs $H=A(\F_p)$ and checks
\[
           \operatorname{ord}(\bar t)=m,\qquad
           \#\langle\bar t,H\rangle=m\,\#H.
\]
Thus the image of $\bar t$ in $R(\F_p)$ has order $m$, proving
$R(\Q)=\langle t\bmod A\rangle$, cyclic of order $m$.
This does not require torsion injectivity at the working prime $2$
for $191$. At $131$, the complement also has $5$-torsion, so norms of
rational and quadratic places construct the actual $A(\F_3)$ rather
than using all $5$-torsion.

Every rational degree-five class lies on one of
$T_i=[5P]+it+A$, $0\le i<m$.
If $T_i\subset\W$, the full coset $i\bar t+H$ must occur.
The displayed indices are exactly those passing the test.
Simplicity of $A$ implies that every positive-dimensional component with
dense rational points is one of these translates.

\subsection{Fixed divisors in the surviving families}
For $g(Y)=1$, indices $1,2,3,4$ are respectively
\[
 \pi^*\Pic^1(Y)+3P,\quad \pi^*\Pic^2(Y)+P,\quad
 \pi^*\Pic^2(Y)+Q,\quad \pi^*\Pic^1(Y)+3Q.
\]
For $g(Y)=2$ only the middle two occur.
These identities follow using $F$ or $2F$ as reference pullbacks.
The line bundles on $Y$ used here are effective.

The displayed divisors are fixed even when $h^0>1$.
For an effective divisor $D$ containing no full fiber,
$\mathcal O_Y$ is a saturated subbundle of $\pi_*\mathcal O_C(D)$
whose quotient has degree $\deg D-e$, where $e=g(C)-2g(Y)+1$.
Saturation follows by intersecting functions from $Y$ with
$\mathcal O_C(D)$ locally; the degree follows by Euler characteristics.
Tensoring with a degree-$d$ line bundle $M$ therefore gives
\[
 H^0(C,\pi^*M(D))=H^0(Y,M)\quad\text{if }d+\deg D<e.
\]
Use $(d,D)=(1,3P),(2,P),(2,Q),(1,3Q)$ as appropriate.
The inequality holds in every row. Since $P,Q$ are distinct points of
one fiber, none of these divisors contains a full fiber.
Every effective representative contains a rational cusp.

\paragraph{The additional cubic family at $83$.}
Let $u^3+4u^2+2u-5=0$. The canonical point
\[
\begin{split}
[\, &2u^2+4u-4,\ 3u^2+5u-7,\ u^2+3u-3,\ 3u^2+3u-5,\\
    &2u,\ -2u^2-2u+2,\ 2\,]
\end{split}
\]
is fixed by $w_{83}$. The polynomial is irreducible and the fifth
coordinate divided by the seventh equals $u$, so the three conjugates
form a rational reduced degree-three divisor $D$.
These exact checks are in \texttt{cubic83.m}.
The points are ramified, so their multiplicity-one divisor contains no
full fiber. The section calculation above shows that $D$ is fixed
throughout $\pi^*\Pic^1(Y)+D$.
On $R$, $w_{83}$ acts by $-1$. For $c=[F+D-5P]\bmod A$,
invariance of $F+D$ gives $2c=5t$.
Since $\#R(\Q)=41$, $c=23t$. This accounts for the extra index.

\paragraph{Functions at $79$ and $103$.}
The enumeration also records all classes of $W^1_5(\F_p)$ above rational
quotient labels. Only labels $2,3$ occur. Any rational special class
therefore lies on $T_2$ or $T_3$ over $\Q$, since the quotient labels
are distinguished on reduction. Both entire translates have fixed cusps.
No degree-five function exists. This test concerns individual special
classes, not merely whole cosets, and is part of
\texttt{quotient\_sieve.m}.

All nine cases now follow from Section~\ref{sec:framework}: every
surviving positive-dimensional translate has a proper rational fixed
divisor, and every function exclusion has been proved.

\section{Route T: one family calculation for fifteen levels}
\label{sec:translates}
Here we force every translate into
$V=B+[F_4]+C\subset\Pic^5(C)$.
Either $\pi:C\to Y$ is a double cover with $g(Y)=2$, with
$B=\pi^*\Jac(Y)$ and $F_4$ twice a fiber, or it is a degree-four map
to a genus-one curve, with $F_4$ a fiber.
Degree-two line bundles on a genus-two curve and degree-one line bundles
on a genus-one curve are effective. Thus $V\subset\W$.
It is a proper image of $B\times C$, of expected dimension $b+1$,
where $b=\dim B$.
\begin{center}\small
\begin{tabular}{@{}rrlrrl@{}}
\toprule $N$&$p$&quotient generators&$b$&$\#B(\F_p)$&subgroups tested\\
\midrule
74&3&$w_{74}$&2&21&all order $7$\\
77&5&$w_{77}$&2&35&all order $7$\\
82&3&$w_2,w_{41}$&1&6&whole $B(\F_p)$\\
86&7&$w_2,w_{43}$&1&8&all cyclic order $8$\\
99&5&$w_9,w_{11}$&1&10&whole $B(\F_p)$\\
111&2&$w_{111}$&2&15&all order $5$\\
118&3&$w_2,w_{59}$&1&5&all order $5$\\
123&2&$w_3,w_{41}$&1&3&all order $3$ or $5$\\
128&3&$w_{128},\,s w_{128}s^{-1}$&1&6&whole $B(\F_p)$\\
141&2&$w_3,w_{47}$&1&3&all order $3$ or $5$\\
142&3&$w_{71}$&2&15&all order $5$ or $7$\\
143&2&$w_{143}$&2&15&all order $3$\\
145&3&$w_5,w_{29}$&1&4&all order $2$\\
155&2&$w_5,w_{31}$&1&3&all order $3$ or $5$\\
159&5&$w_3,w_{53}$&1&6&all order $3$\\
\bottomrule
\end{tabular}
\end{center}
At $128$, $s$ is the rational normalizer action $q\mapsto-q$.
All rows are verified by the single routine \texttt{translates.m}.

\subsection{Why every possible direction is tested}
The possible positive simple factors have these profiles; $1^2$
denotes a repeated elliptic isotypic part, not a simple surface.
\begin{center}
\begin{tabular}{@{}ll@{}}
\toprule Profile&Levels\\
\midrule
$1^2$&$74,86,111,159$\\
$1$&$77,82,99,118,128,143$\\
$1+1$ (nonisogenous)&$123,141,142,155$\\
$1+2$ (nonisogenous)&$145$\\
\bottomrule
\end{tabular}
\end{center}
For prime-order tests, the Frobenius order of every possible simple
constituent is divisible by a tested prime $\ell\ne p$.
Every nonzero $\Q$-direction contains a $\Q$-simple subvariety isogenous
to such a constituent. Isogeny preserves Frobenius polynomials, so its
reduction contains a rational order-$\ell$ point.
The script tests \emph{all} order-$\ell$ subgroups of $J(\F_p)$,
including those from rank-zero factors. Repeated isotypic parts
therefore leave no unsearched directions.

At $82,99,128$, the unique positive elliptic part equals the known $B$.
The identity is checked on canonical differentials using its newform.
The whole $B(\F_p)$ may thus be used, including its $p$-primary part;
this is not a prime-to-$p$ torsion lifting argument.

At $86$, the repeated elliptic factor is isogenous to
$E_{43}:y^2+y=x^3+x^2$. Its $2$-division polynomial is irreducible
over $\Q$, so it has no rational $2$-isogeny.
Every elliptic curve in its $\Q$-isogeny class is linked by an odd-degree
isogeny: remove full multiplication kernels until the kernel is cyclic;
an even cyclic kernel would give a rational order-two subgroup.
At $7$, $E_{43}$ has eight points and a unique nonzero $2$-torsion
point, so its group is cyclic of order eight.
Odd-degree isogenies preserve this $2$-primary group.
Every nonzero possible direction therefore contains a rational cyclic
order-eight subgroup on reduction. All $192$ such subgroups of
$J(\F_7)$ are tested.

\subsection{The certificate and its lift}
Canonical matrices are checked to define the quotient automorphisms
and, in degree four, to commute and generate a group of order four.
Their invariant differential space has dimension $b$.
The order of $B(\F_p)$ is computed from the quotient Jacobian or its
invariant Hecke space. Norms of rational and quadratic places generate
a subgroup of that order, constructing the actual $B(\F_p)$.

For \emph{every} tested $S$ and every whole coset $x+S\subset\W(\F_p)$,
the certificate checks:
\begin{enumerate}
\item $S\subset B(\F_p)$;
\item $x+S\subset B(\F_p)+[F_4]+C(\F_p)$;
\item $\dim T_x Z_S=b+1$;
\item the differential of $B\times C\to\Pic^5(C)$ has rank $b+1$
at each relevant rational preimage.
\end{enumerate}
For the last check, the group-trace kernel on canonical differentials
annihilates $T B$; it is evaluated at every rational point of $C$.
The rank and point-count checks also verify the expected smooth
reduction of this embedded $B$.

The complete surviving calculation is below. Counts are
subgroup--coset pairs, not necessarily disjoint cosets.
\begin{center}\small
\begin{tabular}{@{}rrrrrr@{}}
\toprule $N$&effective divisors&$\#\W(\F_p)$&subgroups&
surviving cosets&$\dim T Z_S$\\
\midrule
74&1340&1316&8&18&3\\
77&12528&12488&1&40&3\\
82&2574&2568&1&6&2\\
86&64384&64384&192&12&2\\
99&8504&8504&1&6&2\\
111&328&320&6&12&3\\
118&3104&3104&6&6&2\\
123&842&842&$1+6$&$6+0$&2\\
128&1812&1812&1&4&2\\
141&566&566&$1+1$&$4+0$&2\\
142&1668&1656&$781+57$&$12+0$&3\\
143&332&324&4&20&3\\
145&1552&1552&15&24&2\\
155&536&536&$4+6$&$4+0$&2\\
159&23476&23476&13&16&2\\
\bottomrule
\end{tabular}
\end{center}
For two prime orders, the two counts follow their order in the preceding
table. All surviving cosets have the displayed tangent dimension.

Let $T\subset\W$ be any positive-dimensional component with dense
rational points. Choose a tested subgroup in its reduced direction.
The checks give a rational point of $B\times C$ above the reduction
of a rational point of $T$. Lift the subgroup as in
Section~\ref{sec:local}. Both $T$ and $V$ lie in $Z_S$.
Lemma~\ref{lem:local} gives $T\subset V$.
At $145$ this treats the elliptic and simple-surface possibilities
together, since both have even Frobenius order at $3$.
No separate surface calculation is needed.

\subsection{Why the known family contributes no quintic points}
Project to $J/B$. The image of $V$ is a translate of the image of $C$,
whose differences generate $J/B$.
In every row $\dim(J/B)>1$. That curve cannot contain a positive-dimensional
abelian translate: such a translate would be an elliptic curve equal to
the image, whose differences generate only a one-dimensional variety.
The image of $T$ is therefore a point, and
$T\subset B+[F_4]+P$ for a fixed geometric point $P\in C$.
Each class there has an effective representative containing $P$.

Every function exclusion has already been proved.
If $T\subset W^1_5$, it contributes no irreducible rational quintic divisor.
Otherwise, on its dense $h^0=1$ open subset, the unique rational
representative contains $P$ and its entire Galois orbit.
This orbit has degree at most four: more than five is impossible,
and degree five would make all classes in the dense open subset equal
to the orbit divisor, contradicting $\dim T>0$.
It is thus a proper rational sub-divisor.
No positive-dimensional component contributes infinitely many quintic
points. Section~\ref{sec:framework} completes the proof.

\section{Reproduction and scope of verification}
This folder is independent of the working repository. The computations
use Magma V2.29-4. From this folder run
\begin{verbatim}
python3 run.py --jobs 8
\end{verbatim}
Up to $32$ independent workers can be requested. To inspect saved output
without recalculation, run
\begin{verbatim}
python3 run.py --check-logs
\end{verbatim}
The runner checks all $36$ certificates for a success marker and error
messages: Magma errors need not produce nonzero exit status.
It writes \texttt{verification.json} with source and log hashes.
Individual scripts run from \texttt{code/}, for example
\begin{verbatim}
magma -b N:=145 translates.m
\end{verbatim}
\begin{center}\small
\begin{tabular}{@{}lll@{}}
\toprule File in \texttt{code/}&Purpose&Runs\\
\midrule
\texttt{arithmetic.m}&Double-cover inequalities&1\\
\texttt{rank\_zero.m}&Rank zero, pencil lifts, $98$ refinement&3\\
\texttt{empty\_cosets.m}&No surviving abelian translates&3\\
\texttt{quotient\_sieve.m}&Rational labels and complete cosets&9\\
\texttt{cubic83.m}&Exact fixed cubic divisor&1\\
\texttt{local\_pencils.m}&Local $W^1_5$ test&3\\
\texttt{pencils74.m}&Two-prime exceptional-label sieve&1\\
\texttt{translates.m}&Known families and tangent spaces&15\\
\bottomrule
\end{tabular}
\end{center}
Helpers \texttt{common.m}, \texttt{coordinates.m}, \texttt{place\_cache.m},
and \texttt{translation.m} supply exact shared operations and coordinate
conversions. The cache accelerates norms; it does not replace exhaustive
degree-five enumeration.

The code checks finite groups, divisor classes, Riemann--Roch dimensions,
automorphisms and tangent ranks. The deductions use the theorems and
proved lemma above; this is not a formal-proof-assistant verification.
The published gonality input is used only for rational tetragonal pencils
at $74,77,85,88$ and gonality six at $82$.
The cusp-field description and degeneracy map at $98$ are standard
modular-curve facts, not numerical searches.

\begin{thebibliography}{9}
\bibitem{DHJO}
M. Derickx, W. Hwang, D. Jeon, and P. Orli\'c,
\emph{Modular curves $X_0(N)$ of density degree $5$},
arXiv:2503.08975v2, Section~7.
\url{https://arxiv.org/html/2503.08975v2}.
\bibitem{NV}
F. Najman and M. Varivoda,
\emph{Torsion groups of elliptic curves that appear infinitely often
over septic, octic and nonic fields},
arXiv:2602.03513v4, Section~5.5, Proposition~5.5.
\url{https://arxiv.org/html/2602.03513v4}.
\bibitem{NO}
F. Najman and P. Orli\'c,
\emph{Gonality of the modular curve $X_0(N)$},
arXiv:2207.11650, Theorems~1.2 and~1.4.
\url{https://arxiv.org/abs/2207.11650}.
\bibitem{KL}
V. A. Kolyvagin and D. Yu. Logachev,
\emph{Finiteness of the Shafarevich--Tate group and the group of
rational points for some modular abelian varieties},
Leningrad Math. J. \textbf{1} (1990), 1229--1253.
\url{https://www.mathnet.ru/aa47}.
\bibitem{Faltings}
G. Faltings, \emph{Diophantine approximation on abelian varieties},
Ann. of Math. \textbf{133} (1991), 549--576.
\url{https://annals.math.princeton.edu/1991/133-3/p04}.
\end{thebibliography}
\end{document}
